% ==================== 结语 ====================
% 本文件只保留创作者结语，不恢复已经删除的技术附录。
\backmatter

\begingroup
\titlespacing*{\chapter}{0pt}{-24pt}{18pt}
\chapter*{结语}
\phantomsection
\addcontentsline{toc}{chapter}{结语}
\markboth{结语}{}

\noindent
这本小册子的正文到这里告一段落。感谢你翻开它，也感谢你读到这里。希望在阅读的过程中，你能有所收获，如果能从中了解到一些 AI 的发展趋势，或者对 AI 的理解有了新的认知，那就再好不过了。

\section*{matto}

还记得 2024 年年底去 CM 时，在整理待买刊物的时候，意外发现了一个名为“篠澤広研究室”的摊位，他们带来的是一本“让篠澤広为大家讲解物理学”的合同志。当时忍不住感慨：喔……原来同人展上还有让小广讲物理学的学术刊物啊。
第二年的夏 CM，学马相关的摊位就百花齐放了。那天一路逛到南馆，看着摊主们忙着布置倾注了各自心意的刊物，不少桌面上还散落着“初星学园創作学園”的彩色传单。也不知道为什么，那一刻眼泪忽然涌了出来。正是在那里，油然产生了一个想法：想自己也做一本属于广的同人刊。
有了想法之后，整个过程都磕磕绊绊的。每天在宇宙厂加班加到怀疑人生，但也正是在工作中对 AI 的大量学习和实际使用，给内容编写提供了很多思路。编写时，也尽可能挑选了近期值得关注的 AI 进展，希望能把它们讲得清楚一些。不过 AI 发展得实在太快，如果各位在阅读时发现不够准确或已经过时的内容，也欢迎来来一起交流~
最后，感谢所有参与这本合刊的朋友。sorodo老大为这本刊物做了很多工作，从 TeX 架构到封面设计，都帮了很大的忙，cinno也是很爽快地回应了征文...没有这些帮助，这个想法估计还会在我的电脑里继续躺上很久很久...


\section*{sorodo}

虽然中间修修改改，现在的小册子和一开始的构想还是有很大的出入的，但是总算也是赶在IFE之前收尾了。
一开始的构想甚至每一篇都想单独加一篇小漫画，但是时间实在不允许和一些事情变故只能忍痛割舍了喵。

甚至这次还是第一次画漫画，不会画QQ人红豆泥私密马赛。
关于Alice小广的封面，其一是参考科研的神秘命名智能体alice和bob；其二是很喜欢爱丽丝症候群老师的插画。
当然，小广自身的造型也很适合alice这一套衣服就是了。

最后，我猜总有封面党不看文本的，被指到的人快去看书上课（其实是画的实在太少了）喵。


\section*{cinno}

See you on the dark side of the moon.

\vfill

\noindent
\begin{minipage}{\linewidth}
  \color{primary!55}\rule{\linewidth}{0.5pt}
  \vspace{0.65em}

  \begin{tabular*}{\linewidth}{@{}>{\bfseries\color{primary}}l@{\hspace{1.2em}}l@{\extracolsep{\fill}}r@{}}
    \multicolumn{2}{@{}l}{\bfseries\color{primary}筱澤広AI小课堂} & \small 2026年8月  \\
    企划 / 主催 & matto & \\
    美术 / 副编 & sorodo & \\
    投稿 & cinno &
  \end{tabular*}
\end{minipage}
\endgroup
